\section{Fase 2}
\setlength{\columnseprule}{0pt}
\onecolumn
\vspace{-0.3cm}
\subsection{Reproducción del Algoritmo PVA en un agente}
\begin{flushleft}
\resizebox{\textwidth}{!}{%
\begin{ganttchart}[
    x unit=0.55cm, y unit title=0.5cm, y unit chart=0.55cm,
    vgrid={*1{dotted}}, hgrid,
    title height=1, title label font=\bfseries\tiny,
    bar/.style={fill=colorH},
    bar incomplete/.style={fill=colorH!25},
    group/.append style={fill=gray!50, draw=gray!50},
    progress label text={},
    bar height=0.7,
    bar top shift=0.15,
    group right shift=0.15, group top shift=0.15, group height=.7,
    milestone/.style={fill=red},
    bar label font=\small,
    bar label node/.append style={anchor=east, align=left, text width=10.5cm},
    group label node/.append style={anchor=east, align=left, text width=10.5cm},
    bar inline label node/.append style={font=\tiny\bfseries\color{white}}]{1}{7}
  \gantttitle{Septiembre 2025}{4}
  \gantttitle{Octubre 2025}{3}\\
  \gantttitle{S1}{1}\gantttitle{S2}{1}\gantttitle{S3}{1}\gantttitle{S4}{1}
  \gantttitle{S1}{1}\gantttitle{S2}{1}\gantttitle{S3}{1}\\
  \ganttgroup[inline=false]{\textbf{F2. PVA agente individual}}{1}{7}\\
  \ganttbar[progress=100,inline=false]{Analizar trabajo seminal del PVA incluyendo citas}{1}{2}
  \ganttbar[progress=100,inline=true]{\color{white}\textbf{100\%}}{1}{2}\\
  \ganttbar[progress=100,inline=false]{Describir la aplicación del PVA en la dinámica del MAS}{2}{3}
  \ganttbar[progress=100,inline=true]{\color{white}\textbf{100\%}}{2}{3}\\
  \ganttbar[progress=100,inline=false]{Describir la principal limitación y condiciones necesarias}{2}{3}
  \ganttbar[progress=100,inline=true]{\color{white}\textbf{100\%}}{2}{3}\\
  \ganttbar[progress=100,inline=false]{Documentar algoritmo a partir del SoA}{3}{4}
  \ganttbar[progress=100,inline=true]{\color{white}\textbf{100\%}}{3}{4}\\
  \ganttbar[progress=100,inline=false]{Implementar en Python y simular entorno estático}{4}{5}
  \ganttbar[progress=100,inline=true]{\color{white}\textbf{100\%}}{4}{5}\\
  \ganttbar[progress=100,inline=false]{Documentar resultados en entornos estáticos}{5}{6}
  \ganttbar[progress=100,inline=true]{\color{white}\textbf{100\%}}{5}{6}\\
  \ganttbar[progress=100,inline=false]{Implementar en Python y simular entorno dinámico}{5}{6}
  \ganttbar[progress=30,inline=true]{\color{white}\textbf{100\%}}{5}{6}\\
  \ganttbar[progress=30,inline=false]{Documentar resultados en entornos dinámicos}{6}{7}
  \ganttbar[progress=30,inline=true]{\textbf{30\%}}{6}{7}\\
\end{ganttchart}
}
\end{flushleft}
